\section{Preliminaries on Kato manifolds}\label{prel}

In the present paper, a Kato manifold will be a compact complex manifold obtained as follows \cite{kato}, \cite{dl84}. We let $\BB:=\{(z_1, \ldots, z_n) \in \CC^n | |z_1|^2+\ldots+|z_n|^2 < 1\}$ be the standard open ball in $\CC^n$. Consider a modification $\pi: \hat{\BB} \to \BB$ at $0$ and a holomorphic embedding $\sigma: \overline{\BB} \hookrightarrow \hat{\BB}$. Glue the two boundaries of $\hat{\BB} \sminus \sigma(\BB)$ via the real analytic CR-diffeomorphism $\gamma:=\sigma \circ \pi$. The result is a compact complex manifold $X(\pi, \sigma)$, named here a \textit{Kato manifold}. Any small neighborhood of the image of $\del\hat\BB$ in $X(\pi,\sigma)$ is a global spherical shell. We shall refer to the couple $(\pi, \sigma)$ as \emph{Kato data} and to $F:=\pi \circ \sigma : (\BB, 0) \to (\BB, 0)$ as \emph{the corresponding germ}.

The class of Kato manifolds introduced in \cite{kato} is slightly more general, as $\pi$ is allowed to be a modification at more than one point, so~that the resulting manifolds are proper modifications of the ones we described above \cite[Lem.\,2.7, Part\,I]{dl84}. However, for the present discussion this generality will not make any difference.

We denote by $q:\wt{X}\to X$ the universal cover of a Kato manifold $X=X(\pi,\sigma)$. The deck group $\Gamma$ of $q$ is canonically isomorphic to $\ZZ$, and we will denote also by $\gamma$ the positive generator of $\Gamma$, which is indeed induced by the map $\sigma\circ\pi$.

We start by settling several general facts that will be needed in the paper.

\begin{lemma}\label{interFB}
Let $F:(\BB,0)\to (\BB,0)$ be a holomorphic germ with $F(\ov\BB)\subset\BB$. Then $\bigcap_{m\in\NN}F^m(\BB)=\{0\}$.
\end{lemma}
\begin{proof}
By hypothesis there exists $0<r<1$ such that $F(\BB)\subset\BB_r$. By the Schwarz lemma \cite[Th.\,6]{shabat}, we have $\|F(z)\|\leq r\|z\|$ for any $z\in\BB$, hence $\|F^m(z)\|\leq r^m$ for any $m>0$ and any $z\in\BB$. This implies that $\lim_{m\to\infty}\sup_{z\in\BB}\|F^m(z)\|=0$, and hence $\mathrm{diam}(\bigcap_{m\in\NN}F^m(\BB))=0$. The conclusion then follows.
\end{proof}

Given a map $F:(\CC^n,0)\to (\CC^n,0)$, define its stable set:
\begin{equation*}\textstyle
W^s(F)=\{z\in\CC^n\mid \lim_{m\to\infty}\|F^m(z)\|=0\}.
\end{equation*}

\begin{lemma}\label{reuniuneFB}
For any holomorphic map $F:(\CC^n,0)\to (\CC^n,0)$ with $F(\ov \BB)\subset\BB$, we~have $W^s(F)=\bigcup_{m\in\ZZ}F^m(\BB)$, where for $m=-k<0$, $F^{-k}(\BB)$ is the preimage of~$\BB$ via $F^k$.
\end{lemma}
\begin{proof}
The previous lemma shows that $\BB\subset W^s(\BB)$. Clearly $F^m(\BB)\subset W^s(F)$ for any $m>0$ and also for any $m<0$. Conversely, if $z\in W^s(F)$ then there exists $m>0$ such that $F^m(z)\in\BB$, which shows the desired equality.
\end{proof}

Suppose that $F:(\CC^n,0)\to(\CC^n,0)$ is given by a Kato data $(\pi,\sigma)$, $F=\pi\circ\sigma$, such that $\pi$ is a modification at $0$.
For any $k>0$, let $H_k:=F^{-k}(0)$ so~that $H_k\subseteq H_{k+1}$ and put $H_\infty:=\bigcup_{k>0} H_k$. Let us set $\Inv(F):=\CC^n\sminus H_\infty$ and $W^*(F):=W^s(F)\cap\Inv(F)$.

\begin{proposition}\label{stableF}
Let $F:(\CC^n,0)\to(\CC^n,0)$ be a holomorphic map with $F(0)=0$ given by a Kato data $(\pi,\sigma)$ via $F=\pi\circ\sigma$, so~that $\pi$ is a modification at $0$ with exceptional divisor $E$. Suppose that $H_\infty=H_m$ for some $m>0$. Let $D$ be the divisor on $X(\pi,\sigma)$ induced by $E$ after gluing. Then we have a biholomorphism $X(\pi,\sigma)\sminus D\cong W^*(F)/\langle F\rangle$.
\end{proposition}
\begin{proof}
If we put $\BB^*:=\BB\cap\Inv(F)$, then by \ref{reuniuneFB} we have $W^*(F)=\bigcup_{m\in\ZZ}F^m(\BB^*)$. Furthermore, by \ref{interFB}, we have:
\begin{equation*}
\bigcup_{m\in\ZZ}\left(F^m(\BB^*)\sminus F^{m+1}(\BB^*)\right)=\bigcup_{m\in\ZZ}F^m(\BB^*)\sminus \bigcap_{m\in\ZZ}F^m(\BB^*)=W^*(F).
\end{equation*}
Now it is clear that $D\cap \bigl(\hat\BB\sminus \sigma(\BB) \bigr)=\pi^{-1}(H_\infty)\cap (\hat\BB\sminus \sigma(\BB))$ and that
\[
\pi:\bigl(\hat\BB\sminus \sigma(\BB)\bigr)\sminus \pi^{-1}(H_\infty)\to \BB^*\sminus F(\BB^*)
\]
establishes a biholomorphism satisfying $\pi\circ\gamma=F\circ \pi$, where $\gamma=\sigma\circ\pi$. This implies in particular that the group $\langle F\rangle$ acts freely and properly on $W^*(F)$ and that $\pi$ induces the desired isomorphism $X(\pi,\sigma)\sminus D\cong W^*(F)/\langle F\rangle$.
\end{proof}

\begin{proposition}\label{embBall}
Let $(\pi:\hat\BB\to\BB,\,\sigma:\BB\to\hat\BB)$ be a Kato data with germ $F=\pi\circ\sigma$ and let $X=X(\pi,\sigma)$ be the corresponding Kato manifold. Then there exists a holomorphic open embedding $\phi:\hat\BB\sminus \{\sigma(0)\}\to \wt{X}$.
\end{proposition}

\Changel
\begin{proof}
We recall that $\wt{X}=\bigsqcup_{m\in\ZZ} W_m \big/\sim$, where $W_m=\hat\BB\sminus \sigma(\BB)$ for any $m\in\ZZ$ and $\sim$ indicates that $W_m$ is glued to $W_{m+1}$ via $\gamma=\sigma\circ\pi$.

Let us put, for any $l\geq 1$, $\wt{X}_l:=\bigsqcup_{1\leq m\leq l} W_m \big/\sim \, \,\subset \wt{X}$. Also, following \cite[\S 1]{iop}, denote by $(\pi_l:\hat\BB^{(l)}\to\BB,\sigma_l:\BB\to\hat\BB^{(l)})$ the Kato data obtained by composing $(\pi,\sigma)$ with itself $l$ times. Then it is easy to see that $\wt{X}_l=\hat\BB^{(l)}\sminus \sigma_l(\BB)$.

If we denote by $\pi_{(l)}:\hat\BB^{(l)}\to\hat\BB^{(l-1)}$ the induced map, then using that $\pi_{(l)}\circ\sigma_l=\sigma_{l-1}\circ F$ we find that $\pi^{-1}_{(l)}$ gives an embedding of $\hat\BB^{(l-1)}\sminus \sigma_{l-1}(F(\BB))$ into $\hat\BB^{(l)}\sminus \sigma_l(\BB)$. Therefore we find inductively, for any $l\geq 1$, open embeddings
\begin{equation*}
\phi_l:Q_l:=\hat\BB\sminus \sigma(F^{l-1}(\BB))\to\hat\BB^{(l)}\sminus \sigma_l(\BB)\subset \wt{X}
\end{equation*}
which satisfy $\phi_{l+1}|_{Q_l}=\phi_l$. Now since $\bigcap_{l\geq 1} F^l(\BB)=\{0\}$ by Lemma~\ref{interFB}, this means that $\bigcup_{l\geq 1} Q_l=\hat\BB\sminus \{\sigma(0)\}$, and thus the family $\{\phi_l\}_{l\geq 1}$ naturally defines an open embedding $\phi:\hat\BB\sminus \{\sigma(0)\}\to \wt{X}$, which concludes the proof.
\end{proof}

